\subsection{Mã hóa bất đối xứng RSA}

Mật mã khóa bất đối xứng (Asymmetric-key Cryptography), hay thuật toán khóa công khai (Public-key Cryptography), là một hệ thống mật mã dựa trên cặp khóa toán học có liên hệ mật thiết: một khóa công khai (Public Key) dùng cho tác vụ mã hóa và một khóa bí mật (Private Key) chuyên biệt cho tác vụ giải mã. RSA (Rivest-Shamir-Adleman) là một trong những giao thức nền tảng thuộc lớp thuật toán này, được chuẩn hóa và triển khai rộng rãi trong các hệ thống truyền thông an toàn.

An toàn cấu trúc của RSA phụ thuộc vào tính bất khả thi về mặt tính toán của bài toán phân tích thừa số nguyên tố (Integer Factorization) đối với các số nguyên lớn. Theo thiết kế, khóa công khai có thể được phân phối công khai trên môi trường mạng mà không gây ảnh hưởng đến tính bảo mật của hệ thống, bởi vì chi phí thời gian để nội suy khóa bí mật từ khóa công khai vượt quá khả năng xử lý của các kiến trúc máy tính truyền thống đương đại.

Hạn chế kỹ thuật của RSA so với thuật toán khóa đối xứng (như AES) nằm ở độ phức tạp tính toán lớn, dẫn đến thông lượng xử lý chậm hơn đáng kể. Vì vậy, trong thực tiễn thiết kế giao thức mạng, RSA không được ứng dụng để mã hóa trực tiếp các luồng dữ liệu tải trọng. Thay vào đó, thuật toán được chỉ định làm cơ sở cho cơ chế bao bọc và phân phối khóa (Key Wrapping \& Distribution).

Trong kiến trúc hệ thống máy chủ (Secure Chat Server), thuật toán RSA với chuẩn khóa 2048-bit (RSA-2048) được triển khai nhằm thực thi quy trình thỏa thuận khóa (Key Exchange Handshake). Cơ chế này đảm nhiệm các chức năng kỹ thuật sau:
\begin{itemize}[label={-}]
    \item Khởi tạo kênh phân phối an toàn: Thực thể máy khách (Client) sử dụng khóa công khai RSA của máy chủ để mã hóa khóa phiên đối xứng AES (Symmetric Session Key) nội bộ. Khóa công khai này được quy định sẵn hoặc trích xuất từ các chứng chỉ phân phối tĩnh.
    \item Ngăn chặn tấn công trung gian (Man-in-the-Middle): Khóa phiên đối xứng ở định dạng bản mã được luân chuyển qua hạ tầng mạng và chỉ cho phép giải mã bởi máy chủ sở hữu khóa bí mật RSA tương ứng. Quá trình này vô hiệu hóa hoàn toàn nguy cơ đánh cắp khóa phiên thông qua kỹ thuật bắt gói tin (Packet Sniffing).
    \item Kiến trúc mã hóa lai (Hybrid Encryption): Sau khi hoàn tất quy trình trao đổi khóa RSA, giao thức kết nối tự động chuyển đổi sang sử dụng AES cho toàn bộ tải trọng thông điệp. Sự tích hợp này giải quyết trọn vẹn bài toán tối ưu hóa tài nguyên tính toán trong khi vẫn duy trì chuẩn mực bảo mật cao nhất đối với khâu xác thực định danh và thiết lập phiên giao dịch.
\end{itemize}
